% !TEX root = main.tex

% ---------------------------------------------------------------------------
% このファイルは単体ではコンパイルできません（\documentclass がない、設定の断片です）。
% 間違えて latexmk preamble.tex としたときに、
% 何千行もエラーを出す前に止めるための仕掛けです。
% ---------------------------------------------------------------------------
\ifdefined\chapter\else
  \errmessage{^^J
   ================================================================^^J
   preamble.tex is a fragment. It cannot be compiled on its own.^^J
   preamble.tex は単体ではコンパイルできません。^^J
   latexmk                 (本編ぜんぶの PDF)^^J
   ================================================================^^J}
  \expandafter\stop
\fi
%% ===========================================================================
%%  パッケージと設定
%%
%%  ここは書き方の設定です。論文を書くのに触る必要はありません。
%%  自分用に書き換えるのは main.tex だけで済むようにしてあります。
%%
%%  余白を変えたい、参考文献の形式を変えたい、といったときだけここを開いてください。
%% ===========================================================================

%% ===========================================================================
%%  参考文献の方式（どちらか一方だけを有効にする）
%% ---------------------------------------------------------------------------
%%   biblatex + biber : .bib をそのまま使え、日本語文献も UTF-8 で素直に通る。
%%   natbib  + bst    : 天文分野で使われてきた aasjournalv7.bst をそのまま使う。
%%  本文の書き方（\citep{} / \citet{}）はどちらでも同じです。
%%  切り替えたら、いちど latexmk -C で中間ファイルを消してから作り直すこと。
%% ===========================================================================
\def\BibStyleBiblatex{}     % ← 既定
%\def\BibStyleNatbib{}      % ← こちらを使うときは上の行をコメントアウトする


%% ===========================================================================
%%  パッケージと設定
%% ===========================================================================

%% 共通プリアンブル（フォント・数式・図・単位・マクロ）。
%% 中身は hh-template.sty を見てください。全テンプレートで共通です。
%% 余白と hyperref はこの下で論文用に設定するので、ここでは外しておく。
\usepackage[nogeometry,nohyperref]{hh-template}

%% 版面。ltjsbook の既定は余白が広いので詰め、両面印刷でのど側を広めに取る。
\usepackage[a4paper,inner=30truemm,outer=22truemm,top=25truemm,bottom=25truemm]{geometry}

%% 図表
%% caption は ltjsbook を知らないので
%%   Package caption Warning: Unknown document class ...
%% という警告を必ず出しますが、実害はありません。無視してください。
\usepackage{caption}
\usepackage{subcaption}      % 図(a)(b) を並べる
\captionsetup{font=small,labelfont=bf}
\captionsetup[sub]{font=small,labelformat=parens,labelfont=bf,%
                   justification=raggedright,singlelinecheck=false}
\usepackage{longtable}       % 何ページにもわたる表
\usepackage{pdflscape}       % 横長の表を 90 度回して 1 ページに収める
\graphicspath{{fig/}}        % \includegraphics{foo} で fig/foo を探す

%% 目次
\usepackage[nottoc,notlot,notlof]{tocbibind}   % 引用文献を目次に載せる
\setcounter{tocdepth}{2}      % 目次に出すのは節まで
\setcounter{secnumdepth}{3}   % 番号を振るのは小節まで

%% 参考文献の実体
\ifdefined\BibStyleBiblatex
  \usepackage[backend=biber,
              style=authoryear,
              natbib=true,        % \citep \citet を使えるようにする
              maxcitenames=2,     % 3 人以上は「Hotta et al.」
              maxbibnames=99,     % 文献リストでは全員書く
              giveninits=true,    % 名は頭文字だけ
              uniquename=init,
              uniquelist=false,   % 3 人以上を必ず「Hotta et al.」に truncate する
              sorting=nyt,
              doi=false,isbn=false,url=false,eprint=false,
              backref=false]{biblatex}
  \addbibresource{thesis.bib}
\else
  \usepackage{natbib}
  \setcitestyle{authoryear,round,aysep={},citesep={;}}
\fi

%% リンク。LuaLaTeX なので日本語のしおりはそのまま通る（pxjahyper は不要）。
%% 紙で提出する版は下の行と入れ替えて文字を黒にすること。
\usepackage[unicode,colorlinks=true,allcolors=blue,pdfencoding=auto]{hyperref}
%\usepackage[unicode,colorlinks=false,pdfencoding=auto]{hyperref}  % 印刷用
\usepackage{url}

%% 表紙
\usepackage{thesiscover}

%% 添削のときだけ使うもの。行番号を振ると「12 行目のここ」と指摘しやすくなる。
%% 事務に提出する版では必ずコメントアウトすること。
%\usepackage{lineno}
%\linenumbers

%% この論文だけで使うマクロが必要ならここに書く。
%% ただし、数式の記法をマクロで置き換えるのはやめてください。
%% ここだけで通じる書き方を覚えても、他所で論文を書くときには使えません。
